\section{Conclusion}
% Length: ~5%
% Focus: Provide a strong, concise wrap-up.

% Must contain:
% 1. Summary: A brief recap of the environmental triggers and internal conflicts.
% 2. Final Takeaway: Reiterate that the RAG environment demands evaluators to evolve from passive scoring models into active investigating agents.
Large language models (LLMs) have emerged as powerful tools for automatic evaluation, offering scalability and flexibility across a wide range of tasks. This survey highlights the evolution of LLM-based evaluation from simple single-pass scoring to more structured and reliable approaches, including rubric-based calibration, multi-agent collaboration, and agentic evaluation frameworks.

Recent methods such as AutoCalibrate~\cite{autocalibrate} and Prometheus 2~\cite{prometheus2} demonstrate that explicit evaluation criteria and specialized evaluator models can significantly improve alignment with human judgments. Similarly, multi-agent approaches like PRD~\cite{prd} and culturally-aware frameworks such as MACD~\cite{macd} show that incorporating diverse perspectives can reduce certain forms of bias. Furthermore, agentic evaluation frameworks~\cite{agent_judge, auto_eval_judge} introduce stepwise verification and task decomposition, enabling more reliable assessment of complex outputs.

Despite these advancements, several critical challenges remain. LLM judges continue to exhibit sensitivity to prompt design and superficial features~\cite{comm_systems}, while cultural bias and value misalignment persist across diverse contexts~\cite{macd}. Multi-agent debate systems, although promising, do not always guarantee genuine reasoning and may suffer from false consensus~\cite{debate_limit}. Additionally, integrating external tools introduces new issues such as tool-memory conflicts~\cite{tool_conflict}. From a systems perspective, advanced evaluation pipelines increase computational cost, latency, and implementation complexity~\cite{agent_survey}.

Taken together, the literature suggests that no single mitigation strategy is sufficient to fully address the limitations of LLM-based evaluation. Instead, the most promising direction lies in hybrid approaches that combine calibrated rubrics, structured reasoning, multi-agent verification, and agentic decomposition, supported by human oversight in high-stakes scenarios.

Future research should focus on improving robustness, fairness, and generalization across tasks and domains, while also addressing practical constraints such as efficiency, scalability, and privacy. Advancing toward trustworthy evaluation will require not only better model architectures, but also principled evaluation frameworks that can adapt to diverse and evolving real-world requirements.